\centerline{\textbf{\large Verificación de consenso e inclusión de transacciones en Ethereum}}
\centerline{\textbf{\large con conocimiento cero aplicado a la transferencia de activos digitales}}

\vspace{2cm}

\rightline{Autor: Javier López - Alejandro Marín}
\vspace{0.5cm}
\rightline{Asesor: PhD Marcos Villagra}

\vspace{0.5cm}

\noindent
\textbf{RESUMEN}

\vspace{0.5cm}

\noindent
Los puentes entre cadenas o popularmente conocidos como \textit{cross-chain bridges}, permiten 
transferir activos entre distintas \textit{blockchains} y constituyen una pieza central de la 
interoperabilidad. La mayoría se apoya en el patrón \textit{lock-and-mint} operado por un 
\textit{relayer} de confianza, un actor que tras observar el bloqueo de fondos en la cadena 
de origen acuña el equivalente en la cadena de destino. Como toda la seguridad recae en ese 
operador, que podría acuñar de forma fraudulenta, los puentes se han vuelto uno de los blancos 
más explotados del ecosistema. 
En este trabajo se diseña e implementa un puente de confianza reducida
(\textit{trust-minimized}) que sustituye esa confianza por verificación criptográfica. El
contrato de destino solo acuña cuando valida \textit{on-chain} una prueba de conocimiento cero
(\textit{zero-knowledge}) que demuestra tanto la inclusión real de la transacción de bloqueo en
un bloque de la cadena de origen a través de una prueba de inclusión en el árbol de transacciones, 
como la coincidencia del destinatario y el monto de la transferencia con esa transacción, de modo que
no puedan adulterarse. La lógica de verificación se desarrolló en Noir y, de forma equivalente,
en Rust para la máquina virtual de conocimiento cero SP1, con verificadores \textit{on-chain} en
Solidity, y ambas implementaciones se comparan empíricamente en rendimiento, costo de
verificación y viabilidad práctica. Un \textit{relayer} \textit{off-chain} construye las pruebas
de consenso y de inclusión de cada operación y las remite al contrato de destino, donde deja de
ser un acuñador de confianza y pasa a ser un probador sin permisos, ya que cualquiera con
pruebas válidas puede liberar los fondos y nadie puede acuñar sin ellas. De manera complementaria se diseña la
validación de la firma del comité de sincronización (\textit{sync committee}) de Ethereum como
raíz de confianza sobre la finalidad del bloque. La solución se validó sobre redes de prueba reales, con Ethereum
Sepolia como cadena de origen y Optimism Sepolia como destino. Adicionalmente incluye una interfaz web y un
modo comparativo que evidencia de forma experimental cómo el esquema clásico admite la
acuñación fraudulenta mientras que el esquema con verificación de conocimiento cero la rechaza.

\vspace{0.5cm}

\noindent
\textbf{Palabras Clave}: puentes entre cadenas (\textit{cross-chain bridges}); pruebas de
conocimiento cero (\textit{zk-SNARKs}); Noir; UltraHonk; SP1; verificación \textit{on-chain};
interoperabilidad \textit{blockchain}; prueba de inclusión (árbol de Merkle-Patricia);
confianza reducida (\textit{trust-minimization}); Ethereum.